% ============================================================
\section{Implementazione}
\label{sec:implementazione}
% ============================================================

Ogni tupla di design viene implementata ed eseguita; le
tabelle CSV non vengono mai modificate. Un fallimento del test osservato a runtime viene
analizzato usando solo i valori restituiti dall'esecuzione del test: se rivela un errore di implementazione del test si corregge il test; se
rivela che l'oracolo assunto in design non regge, la tupla viene riesaminata e la correzione viene scritta in un file CSV
separato per metodo.

\subsection{BrokerImpl}

Il metodo \texttt{initialize}, necessario per costruire un'istanza utilizzabile del SUT, non
è incluso tra i 7 metodi di Category Partition: il proprio Javadoc lo descrive come un passo di setup, non
un'operazione rappresentativa puramente delle funzionalità di BrokerImpl.

Gli argomenti di \texttt{initialize} sono stati impostati basandosi su javadoc e firme; unico input non sufficientemente
documentato è \texttt{managed}. ConnectionRetainModes specificava inoltre i
possibili valori di connMode. I valori sono stati scelti con l'obiettivo di presupporre meno condizioni aggiuntive possibili.
Inoltre l'assenza di diretta documentazione tra i metodi testati e questo metodo suggerisce indipendenza tra loro,
non richiedendo una nuova progettazione dei test.

\subsubsection{isDetached(Object obj, boolean find)}

Per costruire \texttt{obj} è necessario modellare il concetto astratto di \texttt{gestito}, tuttavia la documentazione non
fornisce alcuna informazione a riguardo. In questo caso è stato necessario leggere l'implementazione di \texttt{isDetached},
scoprendo la delega a \texttt{ImplHelper.isManageable(obj)} e l'interfaccia \texttt{PersistenceCapable}.

Il Javadoc di \texttt{isDetached} documenta il comportamento generale della ricerca nel DB ("as a last resort this
method will check whether or not the provided object exists in the DB"), ma non il meccanismo concreto con cui
avviene. Per quello è stato necessario leggere l'implementazione, scoprendo che la ricerca comincia recuperando i
metadata della classe dell'oggetto tramite \texttt{\_repo.getMetaData()}. Su quei metadata viene invocato un metodo
statico di utilità.

La ricerca nel DB non è un meccanismo nuovo/separato, è una chiamata diretta a \texttt{this.find(...)}, già
analizzato come metodo da testare. Siccome \texttt{find()} è chiamato internamente come \texttt{this.find(...)},
l'unico modo di far funzionare tutta la vera logica di \texttt{BrokerImpl} tranne quell'ultima chiamata a
\texttt{find()} che va mockata è uno spy.

Le 10 tuple del CSV diventano 10 test di unità, tutte passate.


\subsection{LoginDialog}

% analisi da svolgere
